% ---- Experiments ------------------------------------------------------------
% This section is largely a TODO until training and evaluation are complete.
% Write the setup and metrics now; fill in results later.

\todo{This section is a template. Fill in all tables and results after
training and evaluation are complete.}

\subsection{Evaluation Metrics}

We evaluate on the held-out test split of 13,896 icons. For each test icon,
the model receives the short VLM caption as the prompt and generates SVG code.
We measure the following:

\begin{itemize}
    \item \textbf{SVG Validity Rate (VR):} fraction of generated outputs that
          parse as well-formed XML with a valid root \texttt{<svg>} element.
          A model that generates syntactically broken SVG is unusable.

    \item \textbf{Rendering Success Rate (RSR):} fraction of valid SVGs that
          render without errors using CairoSVG. Catches semantic issues such
          as malformed \texttt{d} attributes that parse as XML but fail to render.

    \item \textbf{CLIP Score:} cosine similarity between the CLIP \citep{radford2021clip}
          embedding of the rendered icon PNG and the CLIP embedding of the
          text prompt, measuring semantic alignment. We use
          ViT-B/32 \citep{radford2021clip} as the backbone.

    \item \textbf{Mean Path Count (MPC):} average number of \texttt{<path>}
          elements in generated icons. Lower is preferable; professional icons
          typically use 1--5 paths. This measures structural economy.

    \item \textbf{FID (Fréchet Inception Distance):} distributional similarity
          between rendered generated icons and rendered ground-truth icons,
          computed over the test split. Lower is better.
\end{itemize}

\subsection{Baselines}

We compare against two categories of baseline.
\textit{Zero-shot LLM baselines} are general-purpose language models prompted
with the same system and user messages as \method{}, but without any
SVG-specific fine-tuning; the ``(zero-shot)'' label distinguishes these from
models that were explicitly trained for SVG generation.
\textit{Supervised SVG baselines} are models trained end-to-end on SVG
corpora and represent the current open-weight state of the art.

\begin{itemize}
    \item \textbf{Qwen3.5-9B (zero-shot):} the base model without fine-tuning,
          prompted identically to \method{}. Isolates the contribution of
          domain-specific fine-tuning.

    \item \textbf{GPT-5.4 (zero-shot):}
          OpenAI's GPT-5.4 prompted via the API with the same prompts.
          Serves as a strong commercial LLM reference point.
          Parameter count is not publicly disclosed.

    \item \textbf{StarVector-8B} \citep{rodriguez2025starvector}: an
          8B-parameter open-weight model that frames SVG generation as
          code generation from a vision-language backbone; state of the art
          on general SVG vectorization at the time of writing.

    \item \textbf{OmniSVG-8B} \citep{yang2025omnisvg}: 8B-parameter model
          trained on MMSVG-2M including a dedicated icon subset, making it
          the most directly comparable supervised baseline.
\end{itemize}

\subsection{Quantitative Results}

\todo{Fill in Table~\ref{tab:main_results} once evaluation is complete.}

\begin{table}[h]
\centering
\caption{Main results on the \dataset{} test split. Best results in \textbf{bold}.}
\label{tab:main_results}
\begin{tabular}{lrccccc}
\toprule
\textbf{Model} & \textbf{Params} & \textbf{VR (\%)} & \textbf{RSR (\%)} & \textbf{CLIP $\uparrow$} & \textbf{MPC $\downarrow$} & \textbf{FID $\downarrow$} \\
\midrule
\multicolumn{7}{l}{\textit{Commercial zero-shot}} \\
\quad GPT-5.4                        & n/a           & \todo{} & \todo{} & \todo{} & \todo{} & \todo{} \\
\multicolumn{7}{l}{\textit{Open-weight zero-shot}} \\
\quad Qwen3.5-\todo{1.7B or 9B}      & \todo{1.7B or 9B} & \todo{} & \todo{} & \todo{} & \todo{} & \todo{} \\
\multicolumn{7}{l}{\textit{SVG-specific (supervised)}} \\
\quad StarVector-8B                  & 8B            & \todo{} & \todo{} & \todo{} & \todo{} & \todo{} \\
\quad OmniSVG-8B                     & 8B            & \todo{} & \todo{} & \todo{} & \todo{} & \todo{} \\
\midrule
\textbf{\method{} (ours)}            & \todo{1.7B or 9B} & \todo{} & \todo{} & \todo{} & \todo{} & \todo{} \\
\bottomrule
\end{tabular}
\end{table}

\subsection{Ablation Studies}

\todo{Run and fill in Table~\ref{tab:ablation} once training is complete.}

\begin{table}[h]
\centering
\caption{Ablation study on key design decisions.}
\label{tab:ablation}
\begin{tabular}{lcccc}
\toprule
\textbf{Variant} & \textbf{VR (\%)} & \textbf{RSR (\%)} & \textbf{CLIP $\uparrow$} & \textbf{FID $\downarrow$} \\
\midrule
\method{} (full)           & \todo{} & \todo{} & \todo{} & \todo{} \\
\quad w/o curriculum       & \todo{} & \todo{} & \todo{} & \todo{} \\
\quad w/o augmentation     & \todo{} & \todo{} & \todo{} & \todo{} \\
\quad short captions only  & \todo{} & \todo{} & \todo{} & \todo{} \\
\quad 1 epoch (vs.\ 3)     & \todo{} & \todo{} & \todo{} & \todo{} \\
\bottomrule
\end{tabular}
\end{table}
